\begin{answers}{ch:posttraining}{事後学習の全体像とファインチューニング}
\item 平均を取る位置は 13 か所で、全体の平均損失は約 $2.16$ です。
対話は 14 トークンで、最後のトークンには次のトークンがないので、入力の位置は 0〜12 の 13 か所です。
アシスタントの 4 か所の $-\log$ の和は本文の例のとおり $1.073$、残りの 9 か所の和は $9\times3.0=27$ です。
よって平均は $(1.073+27)/13=28.073/13\approx2.16$ です。
マスクありの $0.268$ の約 8 倍で、損失のほとんどが「利用者の質問を当てる」ことに使われてしまいます。
\item $\lambda=0.75$、新しいデータから取る行数は 8 行です。
\eqref{eq:post-mix-loss}の $\lambda$ は元のデータの割合なので、$1-0.25=0.75$ です。
新しいデータの行数は $32\times0.25=8$ 行、残りの 24 行が元のデータです（コードでは \code{n_new = round(32 * 0.25)}）。
\item \code{mix_low} は、忘却が最も少ない代わりに、掛け算の正答率も最も低くなりました（実測。300 ステップ、約 4.5 分）。
結果は、掛け算 7\%、加減算 82\%、漱石の損失 4.881 です（学習前は 0\%、90\%、4.877）。
\tabref{tab:post-cpt}と比べると、漱石の損失はほとんど変わらず、加減算も \code{mix} の 79\% より少しよく保たれました。
一方、掛け算は \code{mix} の 20\%、\code{low_lr} の 11\% より低く、「23*4=」には「16」と答えます。
学習率を下げると出発点から遠くへ動かず、混ぜると新しいデータに使う行が半分になるので、どちらの効果も「新しいことを学ぶ量」を減らす方向に働きます。
忘却を防ぐ強さと、新しい能力の伸びは引き換えの関係にあり、この 300 ステップでは \code{mix} のほうがつり合いがよいと言えます。
\item 正答率は 89.0\% から 83.5\% に下がりましたが、「1 ずらした答え」を出すようにはなりませんでした（実測。200 問中 1 問だけ）。
3,000 件のうち 300 件（10 件に 1 件）の答えを 1 大きくし、\code{run_sft} で本文と同じ 600 ステップ学習しました（200 問の評価なので、数ポイントは揺らぎの範囲です）。
どの問題が書き換えられたかに規則はないので、モデルが学べるのは「答えはたいてい正しく、ときどき 1 大きい」という手本の分布です。
貪欲生成では確率が最大のトークンを選ぶので、多数派の正しい答えが出ます。
それでも、まちがった手本に確率を割いたぶん、正しい答えへの確信が薄れ、正答率が下がりました。
手本の誤りも含めてまねるのが SFT なので、手本の質がそのまま上限になります。
\item 防げます（実測）。
\code{drill_text} で作った「12+35=47」の形の行を 4 行ずつまとめ、全位置を正解にした例を 1,000 件、対話 3,000 件に混ぜて \code{run_sft} で 600 ステップ学習しました。
ドリル形式の正答率は 88.5\% で、学習前の 86.5\% を保ちました（混ぜないと 0\%）。
対話の正答率は 85.5\% で、本文の 89.0\% より少し下がりました（バッチの 4 分の 1 がドリルに回るためと考えられます）。
漱石の損失も 5.92 で、混ぜない場合の 6.33 より小さく抑えられました（学習前は 4.88）。
\secref{sec:post-cpt}の「元のデータを混ぜる」と同じ考え方で、SFT でも忘却を防げます。
\item 見るべき観点は、「足りないのは知識か振る舞いか」と「どのデータが用意できるか」です。
まず Q1 として、関係する文書をプロンプトに入れて答えさせ、プロンプトの工夫だけで足りるかを試します（データは評価用の質問と模範解答の組）。
足りなければ Q2 で、社内文書の生テキスト（数億トークン）で継続事前学習をします。
忘却を防ぐため、元の事前学習データも用意して混ぜます。
次に Q3 で、社内の質問と模範応答の組を作り（本章の計算アシスタントでは 3,000 件でした）、SFT をします（安く済ませるなら LoRA）。
逆の順にすると、SFT で身につけた対話の形が継続事前学習で失われるおそれがあるので、この順番にします。
さらに必要なら Q4 で、応答の良し悪しの比較データを集めて選好最適化をします。
どの段階でも、最後に評価用の質問で正答率を測り、足りなければ上に戻ります。
\item \code{<|assistant|>} と、電卓を呼ぶための目印（\chapref{ch:agent}の \code{<|tool|>}）を学習の対象に含める必要があります。
推論のとき、利用者の発話の後で、次に書く目印をモデル自身に選ばせることになるからです。
どちらの目印を書くかが「答えるか、電卓を呼ぶか」の選択そのものなので、それを正解として学ばせなければ、モデルは選び方を知りません。
\code{encode_chat_example} の \code{learn_role=True} がこの設定です。
一方、利用者の発話と、電卓の実行結果（\code{result}）は、引き続き学習しません。
これらはモデルが書くものではなく、外から与えられるものだからです。
\end{answers}
